% !TEX root = main_memo.tex

\chapter{Precise structure} \label{PreciseStructureFinal}

We have now reached the point where we need to define a finite set of generators, for that we require a third and last operation.

\section{The wedge operation}\label{subsectionWedge}


Finite gluing and pointed gluing allow us to generate all functions with compact domains up to continuous equivalence by \cref{Compactdomains}.
The infinite gluing allows us to build functions which map non converging sequences to non converging sequences.
With the pointed gluing, we can build functions which send converging sequences to converging sequences.

The \emph{wedge operation}, that we now introduce, allows us to construct functions exhibiting two new ``continuous'' behaviours on sequences that appear to be essential for functions with non compact domains.
First, several sequences converging to distinct limits can be mapped to converging sequences that share the same limit.
Second, a non converging sequence can be mapped to a converging sequence.
While gluing and pointed gluing naturally arise from operations on spaces, this new operation is fundamentally an operation on functions, with no obvious counterpart for spaces.  

\index[IdxNotion]{operation!wedge}\index[IdxNotion]{wedge operation}
\begin{definition}
Given a natural number $k\in \N$ and functions $(f_i)_{i\leq k+1}$ in $\functionsonbaire$,
we define a function $f:=\fctwedge (f_0,\ldots,f_{k}\mid f_{k+1})$ called the \emph{wedge} of $(f_0,\ldots,f_{k}\mid f_{k+1})$ as follows.
\index[IdxSymb]{v@{$\fctwedge(f_0,\ldots f_k\mid f_{k+1})$}}

Let $A_i= \pgl \dom f_i$, for $i\leq k$, $A_{k+1+i}=\dom f_{k+1}$ for all $i\in\N$, and $B_j= \gl_{i\leq k+1}\im f_i$ for all $j\in \N$.
The domain of $\fctwedge (f_0,\ldots,f_{k}\mid f_{k+1})$ is $\gl_{i} A_i$ and its range is $\pgl_j B_j$.
Then we set for all $i,j$ in $\N$:
\begin{align*}
f ((i)\conc \iw{0})&=\iw{0}\mbox{ if }i\leq k \\
 f((i)\conc(0)^j\conc(1)\conc x)&=(0)^j\conc(1)\conc(i)\conc f_i(x)\mbox{ if }i\leq k \\
  f((k+1+i)\conc x)&=(0)^i\conc(1)\conc(k+1)\conc f_{k+1}(x)
\end{align*}
\end{definition}

\begin{figure}
\centering
\index[IdxFig]{The Wedge operation}
\begin{tikzpicture} [thick, scale=0.35]
\def \HOne{2.5}
\def \VOne{1.5}
\def \HTwo{5} 
\def \LHOne{8}
\def \LVOne{7}
\def \LVTwo{1}
\def \LHTwo{1}
\def\Angle{0}

\def\Second{20}

\foreach \r in {0,1}
\draw (\Angle:\HOne*\r)-- ++(90+\Angle :\LVOne);

\draw (0,-1) -- (\Angle:\HOne*0);

\foreach \r  in {1,...,4}
\draw (0,-1) ..controls +({0.5\r},0) and +(0,-0.8).. (\Angle:\HOne*\r);

\node at ($({\HOne*5},{\VOne*0})+({\LHTwo},1)$) {$\cdots$};

\foreach \r  in {0,1}
\foreach \s in {0,1,2}
{
\pgfmathparse{20*\s}
\edef\c{\pgfmathresult}
\draw[fill=black!\c!white] ($({\HOne*\r},{\VOne*\s})+({\LHTwo},1)$) ellipse (0.8 and 0.5);   
\draw ($({\HOne*\r},{\VOne*\s})+(0,0.5)$) -- ($({\HOne*\r},{\VOne*\s})+({\LHTwo},0.5)$);

\node at ($({\HOne*\r},{\LVOne})+({\LHTwo},0)$)  {$\vdots$};
}

\foreach \r  in {2,3,4}
{
\pgfmathparse{\r-2}
\edef \s{\pgfmathresult}
\pgfmathparse{20*\s}
\edef\c{\pgfmathresult}
\edef \s{0}

\draw[fill=black!\c!white] ($({\HOne*\r},{\VOne*\s})+({\LHTwo},1)$) ellipse (0.8 and 0.5);   
\draw  (\Angle:\HOne*\r)-- ($({\HOne*\r},{\VOne*\s})+(0,0.5)$) -- ($({\HOne*\r},{\VOne*\s})+({\LHTwo},0.5)$);

}

\def \anglesecond{25}
\draw (\Second,-1)-- ++ (\anglesecond :16); 

\foreach \r  in {0,1,...,2}
\foreach \s in {0,1,...,2}
{
\pgfmathparse{\HTwo*\r}
\edef\x{\pgfmathresult}
\pgfmathparse{180-35*\s}
\edef\y{\pgfmathresult}
\pgfmathparse{20*\r}
\edef\c{\pgfmathresult}
\draw ($(\Second,-1)+(\anglesecond :\x)$)-- ++(\y:2.5);
\draw[fill=black!\c!white] ($(\Second,-1)+(\anglesecond :\x)+(\y:2.5)+(0,0.7)$) ellipse (0.5 and 0.7);
}
\node at ($(\Second,-1)+(\anglesecond :{3*\HTwo})+(110:1)+(0,0.7)$) {$\cdots$};

\draw[->] ($({\HOne*0},{\VOne*2})+({\LHTwo},1)$) ..controls +(4,4) and +(-4,4).. ($(\Second,-1)+(\anglesecond :{2*\HTwo})+(180:2.5)+(0,0.7)$) node[midway,above] {$f_0$};
\draw[->] ($({\HOne*1},{\VOne*1})+({\LHTwo},1)$) ..controls +(4,3) and +(-4,2).. ($(\Second,-1)+(\anglesecond :{\HTwo})+(145:2.5)+(0,0.7)$) node[midway,above] {$f_1$};
\draw[->] ($({\HOne*2},{\VOne*0})+({\LHTwo},1)$) ..controls +(3,2) and +(-4,2).. ($(\Second,-1)+(\anglesecond :0)+(110:2.5)+(0,0.7)$) node[midway,above] {$g$};

\end{tikzpicture}
\caption{The wedge operation $\fctwedge(f_0,f_1\mid g)$.} 
\label{fig:Wedge}
\end{figure}

\begin{remark}\label{rem:infinitewedge}
This finitary operation is closely related to the following infinitary operation.%\rnote{ajouter ici l'exemple que l'on a supprimé avant?}
Given an infinite matrix of functions $f_{i,j}:A_{i,j}\to B_{i,j}$, ${(i,j)\in\N\times \N}$, we can define $\fctwedge(f_{i,j}):\gl_j \pgl_i A_{i,j} \to \pgl_i \gl_j B_{i,j}$ by
\begin{align*}
(j)\conc \iw{0}&\longmapsto\iw{0}\\
 (j)\conc(0)^i\conc(1)\conc x&\longmapsto(0)^i\conc(1)\conc(j)\conc f_{i,j}(x)
\end{align*}
Alternatively, we can first form the function $w=\gl_j \pgl_{i} f_{i,j}$ whose codomain is given by $\gl_j \pgl_{i} B_{i,j}$.
Naturally, each $\pgl_i B_{i,j}$ is a pointed space with basepoint $p_j=\iw{0}$ and their \emph{wedge sum}\footnote{See for instance \cite[p.153]{MR0957919}.} is defined through
the topological quotient $\pi:\gl_j  \pgl_i B_{i,j} \to \bigvee_j \pgl_i B_{i,j}$ with respect to the equivalence relation $\{((j)\conc p_j,(j')\conc p_{j'})\mid j,j'\in \N\}$.
Then, we can define $\fctwedge(f_{i,j})$ by $\pi w$, which explains our choice of name for this operation.


Observe that our finitary operation $\fctwedge (f_0,\ldots,f_{k}\mid f_{k+1})$ can be viewed as the following particular case of the infinitary wedge. Each of first $k+1$ functions $f_0,\ldots, f_k$ are repeated infinitely in each of the first $k+1$ columns, $(f_{i,j})_j=\iw{f_i}$ for $i=0,\ldots, k$, hence we informally refer to them as ``the verticals''.
The function $f_{k+2}$ is repeated infinitely on the diagonal of the remainder of the matrix, say $f_{i,i}=f_{k+2}$ for all $i>k$, so we informally call it ``the diagonal''. All other entries are filled with the empty function.
\end{remark}

Note that for any function $f$ we have $\pgl f \equiv \fctwedge (f| \emptyset)$.

\index[IdxNotion]{Cantor--Bendixson!rank!of a wedge}\index[IdxNotion]{wedge operation!Cantor--Bendixson rank}
\begin{fact}\label{BasicfactsWedge}
Let $k\in\N$, $(f_i)_{i\leq k+1}$ be a sequence in $\functionsonbaire$, and set $f:=\fctwedge(f_0\ldots f_k\mid f_{k+1})$.
\begin{enumerate}
\item The wedge operation preserves countability and Polishness of domain and range, surjectivity, finite-to-oneness, and continuity of functions.
\item If $k>0$, then $f$ is not injective.
\item If $f_i\in\sC$ for all $i\leq k+1$ then so does $f$, and 
\[\CB(f)=\max\big(\set{\CB(f_i)+1}[i\leq k]\cup\set{\CB(f_{k+1})}\big).\]
\end{enumerate}
\end{fact}

As with our other operations on functions, we now to look for upper and lower bound criteria.
We can first adapt the ``upper bound'' sufficient criterion \cref{Pgluingasupperbound} to get an upper bound criterion for the wedge operation.

\index[IdxNotion]{upper bound condition!wedge}
\index[IdxNotion]{wedge operation!upper bound condition}

\begin{proposition}[Wedge as upper bound]\label{Wedgeasupperbound}
%Let $A$ be a topological space, $B$ a 0-dimensional separable metrizable space, $f:A\rao B$ be a continuous function, $k\in\N$,
%and $(f_i)_{i\leq k+1}$ in $\functionsonbaire$.
Let $f:A\to B$ in $\functionsonbaire$ be continuous and $(f_i)_{i\leq k+1}\subseteq \functionsonbaire$ for $k\in \N$.
Suppose that there exist $y\in B$ and a clopen partition $(A_i)_{i\in\N}$ of $A$ with the following properties:
\begin{enumerate}
\item For all $i\leq k$, $(\ray{(f\restr{A_i})}{y,j})_{j\in \N}$ is reducible by pieces to $\iw{f_i}$,
\item $(f\restr{A_i})_{i>k}$ is reducible by pieces to $\iw{f_{k+1}}$,
\item $f(A_i)\rao y$.
\end{enumerate}
Then $f\leq\fctwedge (f_0,\ldots,f_{k}\mid f_{k+1})$.
\end{proposition}
\begin{proof}
Use the hypotheses to fix a family $(I_n)_n$ of pairwise disjoint finite subsets of $\N$ and reductions $(\sigma^D_n,\tau^D_n)$ from $f\restr{A_n}$ to $\gl_{i\in I_n} f_{k+1}$ for all $n>k$,
and to get continuous reductions $(\sigma^V_i,\tau^V_i)$ from $f\restr{A_i}$ to $\pgl f_i$ with $\tau^V_i(\iw{0})=y$ using \cref{Pgluingasupperbound}, for all $i\leq k$. 

Define $\sigma$ as follows: 
\begin{align*}
\sigma(x)=&(i)\conc\sigma^V_i(x) && \text{if $x\in A_i$ for some $i\leq k$,}\\
\sigma(x)=&\sigma^D_i(x) && \text{if $x\in A_i$ for some $i>k$.}
\end{align*}

Define $\tau$ as follows:
\begin{align*}
\tau(\iw{0})=&y\\
\tau((0)^j\conc(1)\conc(i)\conc x)=&\tau^V_i((0)^j\conc(1)\conc x)&&\text{if $i\leq k$ and $j\in \N$,}\\
\tau((0)^{i}\conc(1)\conc(k+1)\conc x)=&\tau^D_{n}((i)\conc x) && \text{if $i\in I_{n}$.}
\end{align*} 
By construction $(\sigma, \tau)$ is a reduction from $f$ to $\fctwedge (f_0,\ldots,f_{k}\mid f_{k+1})$. The function $\sigma$ is continuous as a disjoint union of continuous maps by \cref{lem:ContUnion}.
To see that $\tau$ is continuous, note that $\tau$ is continuous on $U=\dom \tau \setminus \{\iw{0}\}$ again by \cref{lem:ContUnion}.
To use \cref{prop:sufficientcondforcont}, we claim that for every sequence $(x_n)_n$ in $U$ that converges to $\iw{0}$ we have $\tau(x_n)\to y$.
If $(x_n)$ is included in $U_i=\bigcup_j N_{(0)^j\conc(1)\conc(i)}$ for some $i\leq k$, then continuity of $\tau_i^{V}$ ensures that $\tau(x_n)\to \tau^V_i(\iw{0})=y$.
If now $(x_n)_n$ is included in $U_{k+1}=\bigcup_{i} N_{(0)^i\conc(1)\conc(k+1)}$ then $\tau(x_n)\to y$ because we assumed that $f(A_i)\to y$.
Since the sets $U_0,\ldots U_{k},U_{k+1}$ form a finite partition of $U$, the claim follows.
\end{proof}

As a corollary, we obtain that, up to continuous equivalence, any function of the form $\fctwedge (f_0,\ldots,f_{k}\mid g)$ is equivalent to one where the verticals $\set{f_i}[i\leq k]$ form an antichain for continuous reducibility\footnote{We go further in \cref{WedgeReducedForm} and prove the existence of a reduced form for the wedge of continuous functions}. To state this corollary, we use the \emph{domination} order on sets of functions, which will also be used in the sequel. For $F,G$ set of functions we say that $F$ is \emph{dominated} by $G$  if for all $f\in F$ there is $g\in G$ with $f\leq g$. Note that if $G$ is finite, then $G$ is equivalent for domination to any maximal antichain of maximal elements of $G$.
When no confusion is possible, we simply denote by $F\leq G$\index[IdxSymb]{2@{$F\leq G$}} for domination and $f\leq G$\index[IdxSymb]{2@{$f\leq G$}} instead of $\set{f}\leq G$. 


\begin{corollary}\label{cor:wedgeSets}
Assume that $f_0,\ldots,f_k, h_0,\ldots, h_l, g \in \functionsonbaire$ are continuous. If $\set{f_i}[i\leq k]$ is equivalent to $\set{h_j}[j\leq l]$ for domination, then
\[\fctwedge(f_0,\ldots,f_k\mid g)\equiv\fctwedge(h_{0},\ldots,h_{l}\mid g).\]
\end{corollary}

\begin{proof}
To see that the left hand side wedge, denoted by $f$, reduces to right hand side wedge,
for each $i\leq k$ let $\phi(i)=\min \set{j\leq l}[f_i\leq h_j]$ and set $A_j=\bigcup \set{N_{(i)}}[\phi(i)=j]$ for $j\leq l$.
Notice that the sets $A_j$, $j\leq l$ partition $\dom f$ and for $y=\iw{0}$ we have $\ray{(f\restr{A_j})}{y,n}\equiv \gl_{i\in\phi^{-1}(j)} f_i \leq |\phi^{-1}(j)| h_j$ for all $j\leq l$. 
\end{proof}



Finding \enquote{natural} sufficient conditions for when a wedge is a lower bound for a function $f$ has proven to be quite tricky,
in particular to make sure that the map $\tau$ is well-defined. So far, we have not been able to find a condition resembling that of \cref{Pgluingaslowerbound}.
We can however generalize \cref{Pgluingaslowerbound2} which enjoys much stronger hypotheses.
In these specific cases we can \enquote{disjointify} enough pieces of the reduction to make sure that their union is well-defined.


\index[IdxNotion]{lower bound condition!wedge}
\index[IdxNotion]{wedge operation!lower bound condition}

\index[IdxNotion]{Disjointification Lemma}
\begin{lemma}[Disjointification Lemma]\label{DisjointificationLemma}
Let $f:A\to B$ in $\functionsonbaire$ be continuous and $(f_i)_{i\leq k+1}\subseteq \functionsonbaire$ for $k\in \N$.
%Let $A$ and $B$ be 0-dimensional separable metrizable spaces, $f:A\to B$, 
%$k\in\N$, and $(f_i)_{i\leq k+1}$ in $\functionsonbaire$.

Suppose that there exist $y\in\im f$ and $(x_i)_{i\leq k}$ in $f^{-1}(\{y\})$ such that 
\begin{enumerate}
%\item $x_i$ is a continuity point of $f$ for all $i\leq k$,
\item for every $i\leq k$, for every open $U\ni x_i$ there exists $(\sigma,\tau)$ reducing $f_i$ to $f$
with $\im(\sigma)\subseteq U$ and $y\notin \overline{\im(f\sigma)}$, and \label{disjoint_condition1}
\item for every open $V\ni y$, there exists $(\sigma,\tau)$ reducing $f_{k+1}$ to $f$ with $\im(f\sigma)\subseteq V$ and $y\notin\overline{\im(f\sigma)}$.\label{disjoint_condition2}
\end{enumerate}
Then $\fctwedge (f_0,\ldots,f_{k}\mid f_{k+1})\leq f$.
\end{lemma}

\begin{proof} We simultaneously define by induction on $n$, a sequence $(V_n)_n\subseteq B$ of open neighborhoods of $y$ and sequences $(U_{i,n})_n\subseteq A$ of open neighborhoods of $x_i$ for each $i\leq k$ satisfying the following conditions:
\begin{enumerate}[label=\arabic*)]
\item $V_n\rao y$ in $B$,
\item $U_{i,n}\rao x_i$ in $A$, and
\item  for all $n\in\N$, for all $i\leq k+1$, there is a continuous reduction $(\sigma_{i,n},\tau_{i,n})$ from $f_i$ to $f\corestr{V_n\setminus V_{n+1}}$,
satisfying $\im(\sigma_{i,n})\subseteq U_{i,n}$ in case $i\leq k$.
\end{enumerate}

Set $V_0=B$ and $U_{i,0}=A$ for all $i\leq k$.
Fix $n\in\N$ and suppose that all relevant sets are built for $m\leq n$.
For all $i\leq k$, use \cref{disjoint_condition1} on $U_{i,n}\cap f^{-1}(V_n)$ (that is an open set containing $x_i$) to find reductions $(\sigma_i,\tau_i)$ from $f_i$ to $f\corestr{V_n}$.
Use \cref{disjoint_condition2} on $V_n$ to find a reduction $(\sigma_{k+1},\tau_{k+1})$ from $f_{k+1}$ to $f\corestr{V_n}$.
Together, \cref{disjoint_condition1,disjoint_condition2} guarantee that $y$ is not in the closed set $\bigcup_{i\leq k+1}\overline{\im(f\sigma_i)}$, so
we can find a natural number $N$ such that $V_{n+1}:=\nbhd{y}{N}$ is disjoint from $\bigcup_{i\leq k+1}\overline{\im(f\sigma_i)}$.
But by continuity of $f$ at $x_i$, for all $i\leq k$ since $f(x_i)=y\in V_{n+1}$ there exists a natural number $N_i\geq n$ such that $U_{i,n+1}=\nbhd{x_i}{N_i}\subseteq U_{i,n}\cap f^{-1}(V_{n+1})$.
This finishes the inductive construction.

Fix now a bijection $\iota:\N\rao(k+1)\times\N$ and for all $m\in\N$ write $\iota(m)=(\iota_0(m),\iota_1(m))$. For all $m\in\N$,
pick a continuous reduction $(\sigma_{m},\tau_{m})$ from $f_{\iota_0(m)}$ to $f\corestr{V_{m}\setminus V_{m+1}}$ as guaranteed by the above construction.

Name $\varphi$ the inverse of $\iota$ and define $(\sigma,\tau)$ as follows:
\begin{align*}
\sigma((i)\conc\iw{0})&=x_i, \\
\sigma((i)\conc(0)^n\conc(1)\conc x)&=\sigma_{\varphi(i,n)}(x), &&\text{for all $i\leq k$ and all $n\in\N$},\\
\sigma((k+1+n)\conc x)&=\sigma_{\varphi(k+1,n)}(x), &&\text{for all $n\in\N$}.
\end{align*}
Also, we let $\tau(y)=\iw{0}$ and if $z\in \dom \tau_m \subseteq V_m\setminus V_{m+1}$
we set 
\[
\tau(z)=(0)^{\iota_1(m)}\conc (1)\conc (\iota_0(m))\conc \tau_m(z).
\]

We claim that the pair $(\sigma,\tau)$ is a reduction from $g:=\fctwedge (f_0,\ldots,f_k\mid f_{k+1})$ to $f$. By construction we have $g=\tau f\sigma$.
To see that $\sigma$ is continuous, we use \cref{prop:sufficientcondforcont} with $U=\dom g \setminus \set{(i)\conc \iw{0}}[i\leq k]$.
As $\sigma$ is continuous on $U$ by \cref{lem:ContUnion}, we show the second condition of \cref{prop:sufficientcondforcont}.
Take $(x_n)_n$ in $U$ converging to some $(i)\conc \iw{0}$. Therefore there exists $M$ such that for all $n\geq M$ we can write $x_n=(i)\conc(0)^{k_n}\conc(1) \conc x'_n$,
with $k_n\to \infty$. Therefore $\sigma(x_n)=\sigma_{\varphi(i,k_n)}(x'_n)\in U_{i,k_n}$ for all $n\geq M$ and since $k_n\to \infty$ and $U_{i,k_n}\to x_i$ we get $\sigma(x_n)\to x_i=\sigma((i)\conc\iw{0})=\sigma(x)$, as desired.

Continuity of $\tau$ similarly follows from \cref{prop:sufficientcondforcont}, this time applied to the open set $\dom \tau \setminus \set{y}=\bigsqcup_m \dom \tau_m$ by remarking that $\iota_1(m_k)\to \infty$ when $m_k\to \infty$. 
\end{proof}

\section{Finite generation at successors of limits}

As a warm-up for the general case, we can now apply the new operation to complete the analysis of $\sC_{\lambda+1}$ when $\lambda$ is 1 or limit.
This will allow us to see all the ingredients of the proof of the general case in a simplified setting.

We need following elaboration on \cref{InfiniteEmbedOmega}. 

\begin{lemma}\label{InfiniteEmbedOmegaStronger}
Let $n\in \N$ and $X_0,\ldots, X_n$ be infinite subsets of a metrizable space $B$.
Then there are pairwise disjoint infinite sets $Y_i\subseteq X_i$ for all $i\leq n$ such that $\bigcup_{i\leq n} Y_i$ is discrete.
\end{lemma}
\begin{proof}
We prove the statement by induction on $n$. If $n=0$, this is simply \cref{InfiniteEmbedOmega}.
So assume it holds for all collections of size $n$ and let $X_0,\ldots, X_{n+1}$ be infinite subsets of $B$. We distinguish two cases.

First case: there exist distinct $i,j\leq n+1$ such that $X_{i}\cap X_{j}$ is infinite.
Reindexing if necessary, assume $j=n+1$ and define $\tilde{X}_{i}=X_{i}\cap X_{n+1}$ and $\tilde{X}_{k}=X_k$ otherwise.
We can apply the induction hypothesis to $\tilde{X}_0,\ldots \tilde{X}_n$ to obtain pairwise disjoint infinite sets $\tilde{Y}_i\subseteq \tilde{X}_{i}$, for $i\leq n$, whose union is discrete.
To obtain the desired sets for the original collection, it suffices to set $Y_j=\tilde{Y}_j$ for $j\neq i$ and further partition $\tilde{Y}_{i}$ into two infinite sets $Y_i$ and $Y_{n+1}$.

Second case: for all $i,j\leq n+1$ the set $X_i\cap X_j$ is finite. We can make them disjoint: set $\tilde{X}_i=X_i\setminus\bigcup_{j\neq i}X_j$ for all $i \leq n$. 
Apply the induction hypothesis to $\tilde{X}_i$ for $i\leq n$ to get infinite pairwise disjoint sets $Y_i\subseteq \tilde{X}_i$, $i\leq n$ with a discrete union.
Now pick a discrete $Y_{n+1}\subseteq \tilde{X}_{n+1}$ by \cref{InfiniteEmbedOmega} and note that the sets $Y_i$ for $i\leq n+1$ are infinite pairwise disjoint discrete sets. 
If $Y=\bigcup_{i\leq n+1} Y_i$ is discrete as well, we are done.
Otherwise, there exists $i\leq n$ such that either $Y_i\cap \overline{Y_{n+1}}\neq \emptyset$ or $\overline{Y_i}\cap Y_{n+1}\neq \emptyset$.  
Suppose that the former holds, choose $y\in Y_i\cap \overline{Y_{n+1}}$ and an open neighborhood $V$ of $y$ such that $V\cap Y_i= \{y\}$ (since $Y_i$ is discrete).
Note that $Y_{n+1}\cap V$ is infinite and so we can shrink these two sets as follows $Y_{n+1}:=Y_{n+1}\cap V$ and $Y_i:=Y_i\setminus \set{y}$.
Both $Y_i$ and $Y_{n+1}$ are infinite and now $Y_i\cap \overline{Y_{n+1}}=\emptyset$.
The latter case is similar and since $n$ is finite, we can apply this procedure repeatedly until we obtain the desired family $(Y_i)_{i\leq n+1}$.
\end{proof}


As an application, we show that if for infinitely many points $y\in \im f$ we can reduce a single function $g$ to corestrictions of $f$ on arbitrary small neighborhood of $y$, then we have $\omega g \leq f$. We will use the stronger statement which involves a finite set $G$ in place of $G$ in \cref{sectionDoubleSucc}.
Given a (finite) set $G$ of functions in $\functionsonbaire$, we write $\omega G$\index[IdxSymb]{w@{$\omega G$}} for $\omega(\gl G)$.

\begin{lemma}\label{Intertwinereductions}
Let $f:A\to B$ between metrizable spaces and $G\subseteq\functionsonbaire$ be finite.
Suppose that for all $g\in G$ there are infinitely many points $y\in B$ such that for all neighbourhood $V$ of $y$ we have $g\leq f\corestr{V}$.
Then $\omega G \leq f$.
\end{lemma}
\begin{proof}
If $G$ is empty, then $\omega G=\emptyset$ which reduces to any function, so suppose that $G\neq\emptyset$. 
For each $g\in G$, let $X_g$ be the set of points $y\in B$ such that for all neighborhood $V\ni y$ we have $g\leq f\corestr{V}$. %, and set $W=\bigcup_{g\in G}W_g$.
Since the sets $X_g$ are infinite subsets of the metrizable space $B$ and $G$ is finite, we can apply \cref{InfiniteEmbedOmegaStronger} to get an injective map $G\times \N\to B$, $(g,n)\mapsto y^g_n$ with discrete image such that  $y^g_n\in X_g$ for all $g\in G$ and all $n\in \N$. As $B$ is metrizable, we can find pairwise disjoint open sets $V^g_n$ of $B$ such that $y^g_n\in V^g_n$ for all $g\in G$ and $n\in\N$. Therefore $\omega G \leq f$ holds by \cref{Gluingaslowerbound}.
\end{proof}

\begin{lemma}\label{Intertwinereductionsforomegacentered}
%Let $A,B$ be topological spaces \ynote*{ajout! Nécessaire pour \cref{InfiniteEmbedOmega} et pour obtenir des clopens! J'ai enlevé la référence à $U$, car on peut directement appliquer le lemme avec $f\restr{U}$}{with $B$ metrizable $0$-dimen\-sional,} $f:A\to B$, and $G\subseteq\functionsonbaire$ be finite.
%Suppose that for all $g\in G$ there are infinitely many points $y\in B$ such that for all \yerror{this makes sense only if $B$ is $0$-dim!!!}{clopen} $V\ni y$ we have $g\leq f\corestr{V}$.
%Then $\omega G \leq f$.
Let $f\in\functionsonbaire$ be continuous and $G\subseteq\functionsonbaire$ be a finite set of centered functions.
If $\omega g \leq f$ for all $g\in G$, then $\omega G \leq f$. 
Moreover, if $f=\bigsqcup_{i=0}^{n} f_i$ for some $n\in\N$, then
%In particular, \yerror*{Nécessaire pour la preuve!!!}{if $f=\bigsqcup_{i=0}^{n} f_i$ is continuous, then:}
\begin{enumerate}
%\item if $\omega g \leq f$ for all $g\in G$, then $\omega G \leq f$, 
\item if $g$ is centered and $\omega g\leq f$, then $\omega g\leq f_i$ for some $i\leq n$,
\item if $\lambda$ is limit and $\Maximalfct{\lambda}\leq f$, then $\Maximalfct{\lambda} \leq f_i$ for some $i\leq n$.
\end{enumerate}
\end{lemma}
%
%\begin{arnote}{Au sujet de ce lemme}
%Comme rédigé maintenant le lemme ne se lit pas très bien.
%J'ai mis du temps à comprendre le sens du ``In particular", en fait la fonction $f$ de la deuxième partie de l'énoncé n'est pas la même que dans la partie principale.
%A minima l'énoncé de la deuxième devrait être quelque chose du genre
%``Les conditions de la première partie sont satisfaites lorsque la fonction $f$ est continue, on obtient ainsi $\omega G\leq f$.
%Lorsque de plus $f=\bigsqcup_{i\leq n} f_i$ alors on a les deux derniers points."
%
%Par ailleurs, au sujet de la dimension 0 de l'espace $B$, je ne te suis pas totalement.
%Ici aussi on pourrait différencier les deux parties de l'énoncé.
%En fait dans la première partie, après application du \cref{InfiniteEmbedOmegaStronger} on pourrait prendre des ouverts $V^g_n$ deux à deux disjoints (et non pas ouvert-fermés)
%alors la famille $(V^g_n)_{g,n}$ est une partition en ouvert-fermés relatifs. Si les $V^g_n$ ne sont pas fermés dans $B$, étant 2 à 2 disjoints il le deviennent relativement à leur réunion.
%Cela nous permet d'utiliser directement \cref{Gluingaslowerbound} sans devoir passer par \cref{Gluingaslowerbound2} qui nécessite $f\in\functionsonbaire$.
%Je me plante?
%
%Voilà donc ce que je propose: on sépare ce lemme en deux. La première partie avec les hypothèse plus générales comme un premier lemme,
%et la deuxième partie comme un corollaire dans le cas spécifique d'une fonction $f\in\functionsonbaire$ continue.
%Par ailleurs j'ai bien l'impression qu'on utilise uniquement la deuxième partie dans la suite.
%Je peux le faire (ça n'est pas grand'chose) mais je voulais ton assentiment avant.
%Sinon on fait la modif' a minima dont je parlais avant, comme tu préfères!
%\end{arnote}
%
\begin{proof}
Take a center $c_g$ for $g\in G$ and a continuous reduction $(\sigma,\tau)$ from $\omega g$ to $f$. The set $\set{f\sigma((n)\conc c_g)}[n\in\N]$ is infinite and by \cref{Centerinvariance} we have $g\leq f\restr{V}$ whenever $V$ is a neighbourhood of $\sigma((n)\conc c_g)$. So by continuity of $f$, for every $n$ and every neighborhood $U$ of $f\sigma((n)\conc c_g)$ it follows that $g\leq f\corestr{U}$. This implies that $\omega G\leq f$ by \cref{Intertwinereductions}.

For the second part, take a center $c_g$ for $g$ and a continuous reduction $(\sigma,\tau)$ from $\omega g$ to $f$. As $f=\bigsqcup_{i\leq n}f_i$, there exists $i\leq n$ such that $\dom f_i$ contains $\sigma((n)\conc c_g)$ for infinitely many $n$. We conclude as before. Finally for $\lambda$ limit, using \cref{Gluingcohomomorphism} and \cref{JSLgeneralstructure} notice that for every sequence $(\beta_k)_{k\in \N}$ cofinal in $\lambda$, we have $\Maximalfct{\lambda}\equiv \gl_{k\in \N} \Minimalfct{\beta_k}$. Choose an increasing sequence $(\beta_k)_{k\in \N}$ cofinal in $\lambda$ and recall that $\iw{0}$ is a center for each $\Minimalfct{\beta_k}$. If $(\sigma,\tau)$ witness $\gl_{k\in\N} \Minimalfct{\beta_k}\leq \bigsqcup_{i\leq n}f_i$, then there exists $i\leq n$ such that $X=\{k\in\N \mid \sigma((n)\conc \iw{0}) \in \dom f_i\}$ is infinite. Note that the subsequence $(\beta_k)_{k\in X}$ is cofinal in $\lambda$, so $\Maximalfct{\lambda}\equiv \gl_{k\in X} \Minimalfct{\beta_k}\leq f_i$, as before. 
\end{proof}


We are now ready to show that $\sC_{\lambda+1}$ is finitely generated, assuming that $\sC_{<\lambda}$ is \bqo{} under continuous reducibility.  We first isolate the only situation where the wedge operation is necessary.% as it will be used explicitly in the sequel. 

\begin{lemma}\label{Diagonalforlambda+1}
Suppose that $f=\bigsqcup_{n\in\N} f_n:A_n\to B$ for simple functions $f_n\in\sC_{\lambda+1}$ with pairwise distinct distinguished points $y_n$. Assume that the following hold: 
\begin{enumerate}
\item $f_0\equiv \pgl \Maximalfct{\lambda}$,
\item $f_n\leq \Minimalfct{\lambda+1}\glbin \Maximalfct{\lambda}$ whenever $n>0$, and
\item $(y_n)_{n>0}$ converges to $y_0$. 
\end{enumerate}
Then for all clopen neighborhood $U$ of $y_0$ we have $f\leq \fctwedge(\Maximalfct{\lambda}\mid \Minimalfct{\lambda+1})\leq f\corestr{U}$.
\end{lemma}
\begin{proof}
To prove that $\fctwedge(\Maximalfct{\lambda}\mid\Minimalfct{\lambda+1})\leq f\corestr{U}$, we use the Disjointification \cref{DisjointificationLemma} with $y=y_0$. Note that $\pgl\Maximalfct{\lambda}\equiv f_0$ is centered with \cocenter{} $y_0$ and so we can choose a reduction $(\sigma,\tau)$ witnessing $\pgl\Maximalfct{\lambda}\leq f_0\corestr{U}$. We choose $x=\sigma(\iw{0})$. To verify the first condition, let $W$ be any clopen neighborhood of $x$ in $\dom f$, let $W'=W\cap  f^{-1}(U)$. We have $\pgl\Maximalfct{\lambda}\geq f_0\restr{W'}$ and since $x\in W'$ we have $\pgl\Maximalfct{\lambda}\leq f_0\restr{W'}$ by \cref{Centerinvariance}. So picking a witness for the latter reduction, we get the desired reduction $\Maximalfct{\lambda}\leq f_0\corestr{U}$ whose $\sigma$ satisfy $\im \sigma \subseteq W$ and $y_0\notin \overline{\im f\sigma}$ by \cref{Rigidityofthecocenter}. 
For the second condition, let $V$ be a neighborhood of $y_0$ in $U$ and note that there exists $n>0$ with $y_n\in V$, since $(y_n)_{n>0}$ converges to $y_0$. Then choosing a clopen neighborhood $W$ of $y_n$ in $V$, note that $\CB(f\corestr{W})=\lambda+1$ so $\Minimalfct{\lambda+1}\leq f\corestr{W}$ by \cref{Minfunctions}, which grants us with the desired reduction.

Next, we show that $f\leq \fctwedge(\Maximalfct{\lambda}\mid\Minimalfct{\lambda+1})$ using \cref{Wedgeasupperbound}.
Fix a clopen neighborhood basis $(B_k)_k$ at $y_0$ and let $\ray{B}{k}$, $k\in \N$, denote the corresponding clopen rays.
Notice that for $n>0$ we have $y_n\in B_{k_n}$ for some $k_n$ and $f_n\corestr{B_{k_n}} \leq \Minimalfct{\lambda+1} \gl \Maximalfct{\lambda}$. So by \cref{Gluingasupperbound} we can pick a clopen subset $A^1_n$ of $f_n^{-1}(B_{k_n})$ such that $f\restr{A^1_n}\leq\Minimalfct{\lambda+1}$ and $A^0_n=A_n \setminus A^1_n$ satisfy $\CB(f\restr{A^0_n})\leq \lambda$. Note that $f(A^1_n)\subseteq B_{k_n}$ and $k_n\to \infty$ as $n\to \infty$ since $(y_n)_{n>0}$ converges to $y_0$, so $f(A^1_n)$ converges to $y_0$.

Setting $A^0=\bigcup_{n>0}A^0_n$ we obtain a clopen set of $\dom f$ with $f\restr{A^0}\leq\Maximalfct{\lambda}$ by \cref{CBrankofclopenunion}. We define a clopen partition of $\dom f$ by setting $\tilde{A}_0= A_0 \cup A^0$ and $\tilde{A}_n=A^1_n$ for $n>0$. Note that the rays of $f\restr{\tilde{A}_0}$ at $y_0$ all have $\CB$-rank $\leq \lambda$ and so they are all reducible to $\Maximalfct{\lambda}$. Moreover $f\restr{\tilde{A_n}}\leq \Minimalfct{\lambda+1}$ for all $n>0$ and since $f(\tilde{A}_n)_{n>0}$ converges to $y_0$, as desired.
\end{proof}

\index[IdxNotion]{Precise Structure!successor of limit}
\begin{theorem}\label{FGatsuccessoroflimit}
Let $\lambda$ be limit or 1. Suppose that continuous reducibility is \bqo{} on $\sC_{<\lambda}$.
Then $\sC_{\lambda+1}$ is generated by the finite set 
\[\set{\Maximalfct{\lambda},\Minimalfct{\lambda+1},\pgl\Maximalfct{\lambda},\omega\Minimalfct{\lambda+1},
\fctwedge(\Maximalfct{\lambda}\mid\Minimalfct{\lambda+1}),\Maximalfct{\lambda+1}}.\]
\end{theorem}

\begin{proof}
By \cref{finitedegreedamuddafuckaz}, we can assume that $f:A\to B$ in $\sC_{\lambda+1}$ has infinite $\CB$-degree.
By the Decomposition \cref{JSLdecompositionlemma}, $f$ is locally simple of $\CB$-rank $\lambda+1$. {Note that a disjoint union of simple functions of $\CB$-rank $\lambda+1$ with same distinguished point is again simple of $\CB$-rank $\lambda+1$.} So we can write $f=\bigsqcup_{n\in\N}f_n$ with each $f_n:A_n\to B$ simple of $\CB$-rank $\lambda+1$ with distinguished point $y_n$ so that the points $y_n$, $n\in \N$, are pairwise distinct.  

Since $\Minimalfct{\lambda+1}\leq f_n\corestr{V}$ for all $n\in\N$ and neighbourhood $V\ni y_n$ by \cref{Minfunctions}, we obtain $\omega\Minimalfct{\lambda+1}\leq f$ by \cref{Intertwinereductionsforomegacentered}.

We partition $\N=N_0\sqcup N_1$ with $N_1=\set{n\in\N}[f_n\equiv\pgl\Maximalfct{\lambda}]$ and set $Y_i=\set{y_n}[n\in N_i]$, $i=0,1$. We distinguish several cases:
%Consider now $I=\set{i\in \N}[f_i\equiv\pgl\Maximalfct{\lambda}]$.

\begin{descThm}
\item[$N_1$ is infinite] By centeredness of $\pgl\Maximalfct{\lambda}$ we can use \cref{Intertwinereductionsforomegacentered} again and get $\Maximalfct{\lambda+1}\leq f$.
As we know that $f\leq\Maximalfct{\lambda+1}$ by \cref{Maxfunctions}, we obtain $f\equiv\Maximalfct{\lambda+1}$ in this case.

\item[$N_1$ is empty] Then by \cref{simplefunctionslambda+1damuddafuckaz}, for all $i$
we actually have either $f_i\equiv\Minimalfct{\lambda+1}$ or $f_i\equiv\Minimalfct{\lambda+1}\gl\Maximalfct{\lambda}$, but in both cases $f_i\leq\Minimalfct{\lambda+1}\gl\Maximalfct{\lambda}\leq 2 \Minimalfct{\lambda+1}$. Therefore, we obtain using \cref{Gluingasupperbound}:
\[\omega\Minimalfct{\lambda+1}\leq f\leq\gl_if_i\leq\omega(\Minimalfct{\lambda+1}\glbin\Maximalfct{\lambda})\leq\omega(2\Minimalfct{\lambda+1})\equiv\omega\Minimalfct{\lambda+1},\]
so in that case $f\equiv\omega\Minimalfct{\lambda+1}$ and we are also done.

\item[$N_1$ is finite and non empty] Suppose first that for some clopen $U\supseteq Y_1$, the set of $n\in N_0$ such that $y_n\notin U$ is infinite,
then $\omega\Minimalfct{\lambda+1}\leq f\corestr{B\setminus U}$ by \cref{Intertwinereductionsforomegacentered} and $|N_1|\cdot\pgl\Maximalfct{\lambda}\leq f\corestr{U}$ by centeredness and \cref{Gluingaslowerbound}, so
\[f\leq\gl_nf_n\leq (\gl_{n\in N_1}f_{n})\glbin(\gl_{n\in N_0}f_n)\leq|N_1|\cdot\pgl\Maximalfct{\lambda}\glbin\omega\Minimalfct{\lambda+1}\leq f\corestr{U}\glbin f\corestr{B\setminus U}\leq f,\]
and in this case $f\equiv|N_1|\pgl\Maximalfct{\lambda}\gl\omega\Minimalfct{\lambda+1}$, so we are done.

Otherwise, any clopen $U\supseteq Y_1$ contains all but finitely many of the points $y_n$, $n\in N_0$.  
In this case, choose pairwise disjoint clopen sets $U_i\ni y_i$, $i\in N_1$, such that for all $i\in N_1$ the set $P_i=\set{n\in \N }[y_n\in U_i]$ is either infinite or equal to $\set{i}$. 
Since $Y_1$ is included in the clopen set $\bigcup_{i\in I}U_i$, it follows that $R=\N\setminus \bigcup_i P_i$ is finite. By adding $R$ to some infinite $P_i$, we get that $(P_i)_{i\in I}$ is a finite partition of $\N$ and we define $f^{i}=\bigsqcup_{n\in P_i} f_n$.

If $P_i=\set{i}$, then 
\[f^{i}\leq \pgl\Maximalfct{\lambda} \leq f\corestr{U_i},\]
since $f^{i}=f_i\equiv \pgl\Maximalfct{\lambda}$ and $\pgl\Maximalfct{\lambda}$ is centered.

If $P_i$ is infinite, note that the sequence $(y_n)_{n\in P_i\setminus \{i\}}$ converges to $y_i$ and so $f^{i}=\bigsqcup_{n\in P_i} f_n$ falls in the case presented in \cref{Diagonalforlambda+1}. Therefore we obtain $f^{i}\leq \fctwedge(\Maximalfct{\lambda}\mid\Minimalfct{\lambda+1}) \leq f\corestr{U_i}$ in this case.



Now it follows from \cref{Gluingasupperbound} and \cref{Gluingaslowerbound}, that
\[\gl_{i\in I}f\corestr{U_i}\leq f\leq\gl_{i\in I}f^i,\]
so in this last case the function is a finite gluing of $ \fctwedge(\Maximalfct{\lambda}\mid\Minimalfct{\lambda+1})$ and $\pgl\Maximalfct{\lambda}$.
\end{descThm}
\end{proof}


\section{The generators}\label{subsectionGenerators}

We write $ \mathcal{P}^+(F)$\index[IdxSymb]{P@{$\mathcal{P}^+(F)$}} for the set of non-empty subsets of a set $F$.
If $F$ is a set of functions, we define $\pgl\mathcal{P}^+(F):=\set{\pgl F'}[\emptyset\neq F'\subseteq F]$\index[IdxSymb]{ptgl@{$\pgl\mathcal{P}^+(F)$}} and $\omega\set{F}:=\set{\omega f}[f\in F]$.\index[IdxSymb]{w@{$\omega\set{F}$}}
For (finite) sets of functions $F_i$ with $i\leq k+1$ we set $\fctwedge(F_0,\ldots,F_k\mid F_{k+1})=\fctwedge(\gl F_0,\ldots,\gl F_k\mid \gl F_{k+1})$.\index[IdxSymb]{V@{$\fctwedge(F_0,\ldots,F_k\mid F_{k+1})$}}

\index[IdxNotion]{generators}
We now define for every $\alpha<\omega_1$ a set $\centered{\alpha}$ of centered functions and a set $\generator{\alpha}$ of generators. 
We set the base case and the limit cases as follows:
\begin{descThm}
\item[Base case] Let $\centered{0}=\emptyset$, $\generator{0}=\emptyset$, and
$\centered{1}=\set{\Minimalfct{1}}$,
\item[Limit case $\lambda$] Let $\centered{\lambda}=\emptyset$, $\generator{\lambda}=\set{\Maximalfct{\lambda}}$, and
$\centered{\lambda+1}=\set{\Minimalfct{\lambda+1}, \pgl \Maximalfct{\lambda}}$.
\end{descThm}
Then for every $\lambda$ limit or null, we define $\centered{\lambda+n+1}$ and $\generator{\lambda+n}$ by induction on $\N$:
\begin{descThm}
\item[For $n>0$] We let 
\[\centered{\lambda+n+1}=\centered{\lambda+n}\cup \pgl \mathcal{P}^+( \centered{\lambda+n}\cup\omega\set{\centered{\lambda+n}} ),\]
and we define the set of generators at ${\lambda+n}$ by $g\in \generator{\lambda+n}$ if 
\begin{enumerate}
\item $g\in \centered{\lambda+n}$, or 
\item $g\in\omega\set{\centered{\lambda+n}}$, or 
\item $g=\fctwedge(F_0,\ldots, F_k \mid F_{k+1})$ for some distinct $F_0, \ldots F_k \in\mathcal{P}^+(\generator{\lambda+n-1})$ and some $F_{k+1}\subseteq\centered{\lambda+n}$.
\end{enumerate}
We call generators of this last type \emph{wedge generators}.\index[IdxNotion]{generators!wedge generator}\index[IdxSymb]{Calpha@{$\centered{\alpha}$}}\index[IdxSymb]{G@{$\generator{\alpha}$}}
\end{descThm}

Note that by our convention $\generator{0}=\emptyset$ generates the only function of null $\CB$-rank, namely the empty function.

\begin{fact}\label{AlreadyKnownGenerators}
\begin{enumerate}
\item We have $\generator{1}=\{\Minimalfct{1}, \omega \Minimalfct{1}\}$,
\item For $\lambda$ limit, the set $\generator{\lambda+1}$ is equivalent (for domination) to 
\[\set{\Maximalfct{\lambda},\Minimalfct{\lambda+1},\pgl\Maximalfct{\lambda},\omega\Minimalfct{\lambda+1},
\fctwedge(\Maximalfct{\lambda}\mid\Minimalfct{\lambda+1}),\Maximalfct{\lambda+1}}.\]
\end{enumerate}
\end{fact}
\begin{proof}
Note that $\centered{0}=\emptyset$, $\generator{0}=\emptyset$, and $\centered{1}=\{\Minimalfct{1}\}$, so $\generator{1}=\{\Minimalfct{1},\omega \Minimalfct{1}\}$.
For the second item, if $\lambda<\omega_1$ is limit, we have $\centered{\lambda+1}=\set{\Minimalfct{\lambda+1},\pgl\Maximalfct{\lambda}}$ and note that if $\pgl \Maximalfct{\lambda}\in F$ then $\fctwedge(\Maximalfct{\lambda}\mid F)\equiv\Maximalfct{\lambda+1}$.
%\ynote*{Do we need the case $\lambda=1$?}{If $\lambda=1$, then $\centered{2}=\set{\Minimalfct{1}, \pgl \Minimalfct{1}, \pgl \omega \Minimalfct{1}}=\set{\Minimalfct{1},\Minimalfct{2},\pgl \Maximalfct{1}}$. Recall that $\Maximalfct{1}=\omega \Minimalfct{1}$. Note that $\fctwedge(\Maximalfct{1}\mid F)\equiv \fctwedge(\Maximalfct{1}\mid F\setminus\{\Minimalfct{1}\})$ and similarly $\fctwedge(\Maximalfct{1}\mid F)\equiv \Maximalfct{2}$ whenever $\pgl \Maximalfct{1}\in F$. Finally note that $\Minimalfct{1}$ is not needed, since $f\glbin \Minimalfct{1}\equiv f$ for all scattered functions with infinite range (more detailed needed).}
\end{proof}


Here are a few useful facts on the combinatorics of generators.
The notation $\sC_{[\lambda,\alpha]}$ stands for $\sC_{\leq \alpha}\cap\sC_{\geq\lambda}$

\begin{proposition}\label{BasicsOnGenerators}
Let $\alpha=\lambda+n$ be a countable ordinal with limit or null and $n\in\N$.
\begin{enumerate}
\item We have $\centered{\lambda+n}\subseteq\centered{\lambda+n+1}$ and $\generator{\lambda+n}\subseteq\generator{\lambda+n+1}$.
\item The sets $\centered{\alpha}$ and $\generator{\alpha}$ are included in $\sC_{[\lambda,\alpha]}$.
\item The sets $\centered{\alpha}$ and $\generator{\alpha}$ are finite. 
\item For all wedge generator $f\in\generator{\alpha}$ there are $F,H\subseteq\centered{\alpha}$ such that $\{\omega H\}\cup F\leq f$ and $f\leq(\gl F)\glbin \omega H$. \label{BasicsOnGenerators_boundingwedge}
\item For all $g\in \generator{\alpha}$, $\omega g$ is equivalent to $\omega H$ for some $H\subseteq \centered{\alpha}$.
\item If $f_n\in\FinGl{\generator{\alpha}}$ for all $n\in\N$, then $\gl_{n\in\N}f_n$ is equivalent to a function in $\FinGl{\generator{\alpha}}$.\label{trickforfinalproof}
\end{enumerate}
\end{proposition}


\begin{proof}%\Needspace*{2\baselineskip}
\begin{enumerate}
\item By induction, using the definition of the sets $\centered{\alpha}$ and $\generator{\alpha}$.
\item Use \cref{CBrankofPgluingofregularsequence1} and \cref{BasicfactsWedge}.
\item It follows from the finiteness of the power set of a finite set and the fact that, in the definition of $\generator{\alpha}$, the requirement $F_i\neq F_j$
for wedge generators forces the sequence of sets $F_i$ to be finite, thus yielding finiteness of $\generator{\alpha}$.
\item Take $g=\fctwedge(F_0,\ldots, F_k \mid F_{k+1})$ for $k\in\N$ with $F_i\in\mathcal{P}^+(\generator{\alpha-1})$ for $i\leq k$ and $F_{k+1}\subseteq\centered{\alpha}$.
Set $F=\set{\pgl F_i}[i\leq k]$ and $H=F_{k+1}$. We have $F\leq g$, $\omega H\leq g$ and $g\leq\gl F\glbin \omega H$ by \cref{Gluingasupperbound_cor}.

\item In case $g$ is a wedge generator, take $F,H\subseteq \centered{\alpha}$ given by the previous point. By \cref{Gluingasupperbound,Gluingaslowerbound2} we get $\omega g\equiv \omega (F\cup H)$.

\item Up to continuous equivalence we can suppose that all $f_n$ are actually in $\generator{\alpha}$.
Note that by finiteness of $\generator{\alpha}$ there is $N\in\N$ and a partition $I_0,\ldots,I_N$ of $\N$ such that for all $n\leq N$ and $i,j\in I_n$
we have $f_i=f_j=g_n$ for some $g_n\in\generator{\alpha}$.
Now if $I_n$ is infinite, we have $\gl_{i\in I_n}f_i\equiv \omega g_n$, which is equivalent to a finite gluing of functions in $\omega \{\centered{\alpha}\}$ by the previous point.
\qedhere
\end{enumerate}
\end{proof}


Our goal is to prove:

\index[IdxNotion]{Precise Structure!Theorem}
\begin{theorem}[Precise Structure]\label{PreciseStructureThm}
For all $\alpha<\omega_1$, every function in $\sC_\alpha$ is continuously equivalent to a finite gluing of functions in $\generator{\alpha}$.
\end{theorem}

For $\alpha<\omega_1$, we call \emph{finite generation at level $\alpha$} the statement:
\begin{descThm}
\item[$\FG(\alpha)$] Every function in $\sC_\alpha$ is equivalent to a finite gluing of functions in $\generator{\alpha}$.\index[IdxSymb]{FG@{$\FG(\alpha)$, $\FG(<\alpha)$, $\FG(\leq\alpha)$}}
\end{descThm}
The notations $\FG(<\alpha)$ and $\FG(\leq\alpha)$ stand for the conjunction of $\FG(\beta)$ for all $\beta<\alpha$ and for all $\beta\leq\alpha$, respectively.
We close this section by stating the results that we have already established about the statements $\FG(\alpha)$.


\index[IdxNotion]{finitely generated set of functions}
\begin{proposition}\label{FGconsequences} We have $\FG(0)$, $\FG(1)$ and $\FG(\lambda)$ for every $\lambda$ limit. 
Moreover for every countable ordinal $\alpha=\lambda+n$, with $\lambda$ limit or null and $n<\omega$, we have
\begin{enumerate}
\item $\FG(<\lambda)$ implies $\FG(\lambda+1)$.
\item $\FG(\leq\alpha)$ implies that every function in $\sC_{[\lambda,\alpha]}$ is a finite gluing of functions in $\generator{\alpha}$.
\item $\FG(<\alpha)$ implies that continuous reducibility is a \bqo{} on $\sC_{< \alpha}$.
\item If $\FG(<\alpha)$ holds then every function in $\sC_\alpha$ is locally centered.
\item If $\FG(<\alpha)$ holds then every centered function in $\sC_{[\lambda,\alpha]}$ is equivalent to some function in $\centered{\alpha}$. In particular, it is equivalent either to $\Minimalfct{\lambda+1}$ or to $\pgl G$ for some $G\subseteq \centered{\alpha-1}\cup \omega\set{\centered{\alpha-1}}$.  
\end{enumerate}
\end{proposition}
\begin{proof}
Note that \cref{JSLgeneralstructure} trivially implies $\FG(\lambda)$ for $\lambda$ limit or null. Using \cref{AlreadyKnownGenerators}, $\FG(1)$ amounts to \cref{LocallyConstantFunctions}. 
\begin{enumerate}
\item Note that $\FG(< \lambda)$ implies that $\sC_{<\lambda}$ is \bqo{} by \cref{FGgivesBQO_2}, and so $\FG(\lambda+1)$ follows from \cref{FGatsuccessoroflimit,AlreadyKnownGenerators}. 
\item Recall that $\generator{\lambda+k}\subseteq \generator{\lambda+k+1}$ for all $k<\omega$, so $G_{\alpha}=\bigcup_{k\leq n} G_{\lambda+k}$.
\item This follows from \cref{SecondstepforBQOthm,FGgivesBQO_2}, 
\item This follows from the previous item combined with \cref{LocalCenterednessFromBQO}.
\item Observe that if $\alpha$ is limit, there are no centered functions in $\sC_\alpha$ and so the statement is vacuously true.
So suppose that $\alpha$ is a successor ordinal.
Since by $\FG(<\alpha)$ the set $\generator{\alpha-1}$ finitely generates $\sC_{[\lambda,\alpha-1]}$, by \cref{finitenessofcenteredfunctions} any centered function $f\in\sC_{[\lambda,\alpha]}$ is
either $f\equiv \Minimalfct{\lambda+1}\in \generator{\alpha}$ or $f=\pgl F$ for some nonempty finite set $F\subseteq\generator{\alpha-1}$.
We show that if $g\in \generator{\alpha-1}$ is a wedge generator, then $\pgl g$ is equivalent to some centered generator in $\centered{\alpha}$. Use \cref{BasicsOnGenerators}~\cref{BasicsOnGenerators_boundingwedge} to get finite sets $G, H\subseteq\centered{\alpha-1}$ with $ \set{\omega H}\cup G \leq g$ and $g\leq \omega H \glbin (\gl G)$.
Note that both $G$ and $H$ are included in $\centered{\alpha-1}$ and that $\iw{g}$ is equivalent by pieces to $\iw{F_g}$ with $F_g=\iw{G\cup\set{\omega h}[h\in H]}$. Hence $\pgl g$ is equivalent to some function $h\in\centered{\alpha}$ by \cref{Pgluingasupperbound}. 
Therefore, forming a set $F'$ where each wedge generator $g$ in $F\subseteq\generator{\alpha-1}$ is replaced by the functions in $F_g$, it follows that $\pgl F\equiv \pgl F'$ by \cref{Pgluingasupperbound},
as desired.    \qedhere
\end{enumerate}
\end{proof}


